\chapter*{序言}
\addcontentsline{toc}{chapter}{序言}

这本书是整个系列的第一册。它的目标不是让读者“看懂一些汇编”，而是建立一套能继续学习系统、编译器、体系结构和性能工程的基础模型：

\begin{center}
\begin{minipage}{0.82\linewidth}
\begin{verbatim}
C/C++ 源码语义
  -> 编译、汇编、链接
  -> x86-64 汇编
  -> 指令、寄存器和内存
  -> ABI、调用约定和二进制接口
  -> 工具观察和可信性能测量
\end{verbatim}
\end{minipage}
\end{center}

本书面向刚学完 C 语言和 C++ 基础、希望进入现代计算机体系结构、操作系统、编译器和高性能计算的读者。它会从“程序如何运行”“CPU 如何执行指令”“内存中的字节如何解释”这些基础问题讲起，再进入 x86-64 汇编、System V AMD64 ABI、C++ 与汇编互操作，以及性能实验的基本方法。

第二册再展开现代 CPU 微架构、cache/TLB、SIMD、HPC kernel、GEMM 和 AI CPU 算子优化。第一册先把地基打稳：读者必须能从一段 C/C++ 代码追到汇编和运行时证据，能解释寄存器、内存、栈、调用约定和测量误差，才能进入后续的优化主题。

本书写作遵循一个原则：任何性能结论都必须能被解释、观察和验证。只看源码不够，只看汇编不够，只看运行时间也不够。真正可靠的性能工程需要把抽象、实现、工具和实验连接起来。

\section*{本书不是速查手册}

本书不是指令表，不是编译参数列表，也不是算法模板合集。每个核心机制都会按以下问题展开：

\begin{itemize}
  \item 它为什么存在？
  \item 如果没有它，会遇到什么问题？
  \item 它给程序员提供了什么抽象？
  \item 底层大致如何实现？
  \item C/C++ 代码如何映射到这个机制？
  \item 如何用工具观察它？
  \item 它对性能有什么影响？
  \item 初学者最容易误解什么？
\end{itemize}

\section*{学习要求}

读者需要已经写过基本 C/C++ 程序，理解变量、函数、指针、数组、结构体、类、模板的基础语法。除此之外，本书会补齐学习汇编、CPU 执行模型、二进制接口和性能测量所需的底层知识。

每一章都包含实验和作业。只读正文不算完成。你需要运行命令、观察输出、写代码、看汇编、做 benchmark，并写报告。
